% Input asset for encounter_multimodal_prr_v2.tex
\begin{table}[!t]
\caption{Finite-window design recipe implied by Theorem~1.  Each step
cites the statement that licenses it; the recipe prescribes modality only
inside the declared window $I$.
\label{tab:design-recipe}}
\begin{ruledtabular}
\begin{tabular}{p{0.94\columnwidth}}
\textbf{Given} finite $d\ge2$ and $m\ge1$, a window
$I=[\tau,T]\subset(0,\infty)$, and target times
$\tau<t_1<\cdots<t_m<T$:\\[3pt]
\textbf{(1) Transport.}  Choose $(\gamma,z_0,\bar z)$ with
$z_0\ne\bar z$ so that the midpoint mean
$\mu(t)=\bar z+(z_0-\bar z)\e^{-\gamma t}$ traverses the design range
monotonically with $x'(t)$ bounded away from zero on $I$
[Eq.~\eqref{eq:exact-m-centres}].  On a compact window $x'>0$ holds
automatically; the quantitative point is to keep
$T\lesssim{}$a few $\gamma^{-1}$ so $x'(t_m)$ is not exponentially
small (late peaks widen as $\sigma/x'(t_j)$); at the $m=3$ anchors these
temporal widths are $0.07,0.15,0.50$ at $\varepsilon=0.1$ and
$0.14,0.30,1.01$ at $\varepsilon=0.2$, the last exceeding the remaining
window $T-t_3=0.7$, so choose $T-t_m$ larger than the last width to keep the
final peak inside $I$.\\[3pt]
\textbf{(2) Support.}  Place the slab centres at $c_j=x(t_j)$,
i.e.\ physical centres $\mu(t_j)$
[Eq.~\eqref{eq:exact-m-centres}].\\[3pt]
\textbf{(3) Width.}  Choose $\varepsilon$ (and $\rho$) so that
$\sigma$ of Eq.~\eqref{eq:exact-m-mixture} obeys
$\min_j(c_{j+1}-c_j)\ge C\sigma$ with $C\approx2.3$.  For near-equal
weights the crossover analysis
[Eqs.~\eqref{eq:exact-m-crossover}--\eqref{eq:exact-m-valley}] then yields
one valley per adjacent gap with nonadjacent components exponentially
small; unbalanced weights shift the crossover by
$\sigma^2\log(w_j/w_{j+1})/(c_{j+1}-c_j)$ and can remove the valley
altogether (at $2.3\sigma$ a two-slab mixture is already unimodal for
weights $0.56/0.44$), so reduce $\varepsilon$ further for a wide
allocation family and rely on step (4).
Theorem~1 guarantees only that some $\varepsilon<\varepsilon_0$ works;
$C\approx2.3$ is a practical proxy for near-equal weights, slightly above the equal-weight
two-Gaussian bimodality threshold $2\sigma$.\\[3pt]
\textbf{(4) Verify the free clock.}  Evaluate
$G_{\varepsilon,\bm w}(t)=a_\varepsilon(t)H_{\sigma,\bm w}(x(t))$
[Eqs.~\eqref{eq:exact-m-factorization}--\eqref{eq:exact-m-mixture}]
on $I$---a cheap one-dimensional computation; the contact factor
$c_{d,\varepsilon}(t)$ follows from Gaussian and wrapped-Gaussian
marginals---and confirm exactly $m$ nondegenerate maxima and $m-1$
nondegenerate minima, with clear root separation, curvature signs, and
endpoint slopes.\\[3pt]
\textbf{(5) Budget.}  Theorem~1 guarantees the complete signature for
$B<B_{\rm top}(\varepsilon;\mathcal W_{w_*})$, where $B_{\rm top}$ is the
topology-retention threshold uniform over the declared allocation
family; equivalently, $f_{B,\varepsilon,\bm w}/B$ inherits the complete
stationary signature
of $G_{\varepsilon,\bm w}$ [Theorem~1;
Eqs.~\eqref{eq:mixed-jet},~\eqref{eq:exact-m-weak-budget}].  The
theorem asserts that this threshold is positive.  For one fixed allocation,
the Supplemental Material gives a sufficient expression
$B_{\rm cert}(\varepsilon;\bm w,\{U_i\})
\le B_{\rm top}(\varepsilon;\{\bm w\})$; its displayed decimal evaluations
are high-precision floating-point values, not interval certificates.  At
practical scales, use the operational classifier crossing
$B_{\rm op}(\varepsilon;m,\bm w,\mathcal P)$ only as a separate numerical
gate and re-verify $f$ directly at the chosen $B$
(Sec.~\ref{sec:finite-scope}).\\[3pt]
\textit{Caveat (quantifier order).}  The choices are sequential: the
geometry and admissible scales may depend on $(d,m)$; one first fixes
$\varepsilon<\varepsilon_0$ and only then
$B<B_{\rm top}(\varepsilon;\mathcal W_{w_*})$.  Nothing
is claimed about modality outside $I$, about an event-mass floor, or
about uniformity of $\varepsilon_0$ or $B_{\rm top}$ in $d$ or $m$
(Secs.~\ref{sec:finite-scope} and~\ref{sec:discussion}).\\
\end{tabular}
\end{ruledtabular}
\end{table}
